\subsection{nml\_numerics}
\label{SEC:nml_numerics}
\inp*[nmlNumerics]{}[nml\_numerics][]

This \hyperref[SEC:User_Input]{namelist block}
is used to specify parameters for the numerical method.

\begin{nmlblocklst}[
,label={LST:nml_numerics}
,caption={Namelist nml\_numerics.}
]
&nml_numerics
  All_Orders                      = -3
  Solution_Order                  =  3
  Projection_Order                = -3
  Quadrature_Order                = -3
  Error_Order                     = -3
  All_Points                      = -1
  Loc_Solution_Pts                =  1
  Loc_Flux_Pts                    =  1
  Loc_Quadrature_Pts              =  1
  Loc_Error_Pts                   =  1
  Loc_Triangle_Pts                =  1
  Loc_Tetrahedra_Pts              =  1
  Invs_Flux_Method                =  1
  Visc_Flux_Method                =  2
  Visc_Variables_To_Ave           =  1
  Combine_Common_Flux_Projections =  0
  Correction_Function             =  1
/
\end{nmlblocklst}

%==========================================================================

\begin{nmldescitem}
All_Orders = -3
\end{nmldescitem}
\inp*[nmlNumerics::AllOrders]{}[All\_Orders][]

\begin{nmldescbody}
If $>0$, sets
\inp[nmlNumerics::SolutionOrder]{Solution\_Order},
\inp[nmlNumerics::ProjectionOrder]{Projection\_Order}, and
\inp[nmlNumerics::QuadratureOrder]{Quadrature\_Order}
all to this value.

\NOTE
This overrides the values given for these three input parameters.
\end{nmldescbody}

%==========================================================================

\begin{nmldescitem}
Solution_Order = 3
\end{nmldescitem}
\inp*[nmlNumerics::SolutionOrder]{}[Solution\_Order][]

\begin{nmldescbody}
Specifies the degree of the polynomial space, \p{S}, to use for the 
solution. 
% 
The order-of-accuracy is approximately \p[+1]{S}.

\NOTE
If $\inp[nmlNumerics::AllOrders]{All\_Orders} > 0$, then
that value will override this input.
\end{nmldescbody}

%==========================================================================

\begin{nmldescitem}
Projection_Order = -3
\end{nmldescitem}
\inp*[nmlNumerics::ProjectionOrder]{}[Projection\_Order][]

\begin{nmldescbody}
Specifies the degree of the polynomial space, \p{P}, to project the residual
onto when using over-integration by residual projection.
%
The residual will be computed as a polynomial of degree \p{S} and projected down
to a polynomial of degree \p{P}. This can potentially improve stability if
aliasing errors are suspected of destabilizing a simulation. 
%
If $\inptext{Projection\_Order} \leq 0$,
this will be set to the same value as \inptext{Solution\_Order}.

\NOTE
If $\inp[nmlNumerics::AllOrders]{All\_Orders} > 0$, then
that value will override this input.
 
\NOTE
This is only used when
$0<\inptext{Projection\_Order}<\inptext{Solution\_Order}$,
and using this means the order-of-accuracy of the simulation is
now approximately \p[+1]{P} instead of \p[+1]{S}.
\end{nmldescbody}

%==========================================================================

\FINALFIX{\newpage}
\begin{nmldescitem}
Quadrature_Order = -3
\end{nmldescitem}
\inp*[nmlNumerics::QuadratureOrder]{}[Quadrature\_Order][]

\begin{nmldescbody}
Specifies the degree of the quadrature polynomial space, \p{Q}, to use when
using over-integration by over-sampling at the quadrature points.
%
If $\inptext{Quadrature\_Order} \leq 0$,
this will be set to the same value as \inptext{Solution\_Order}.

Certain operations during a given residual evaluation
(\forexample evaluating the common\slash upwind interface fluxes)
produce either polynomials that are not within the solution polynomial space
(\forexample of degree greater than the solution order)
or functions that are not polynomials at all (\forexample pressure).
%
By setting
  $\inptext{Quadrature\_Order} > \inptext{Solution\_Order}$,
these functions will be over-sampled at the quadrature points to produce
approximations of these functions as polynomials of degree \p{Q},
which are in turn projected down to polynomials of degree \p{S}
before adding their respective contribution to the residual.
%
This polynomial representation of degree \p{S} is generally more
accurate than directly evaluating the function at the
\p{S} polynomial points, and
can potentially improve stability if aliasing errors are suspected
of destabilizing a simulation.
%
See \sref{SEC:Aliasing} for an example of how \key{aliasing error}[=] arises
and the process of using over-integration by over-sampling at the
quadrature points.

\NOTE
If $\inp[nmlNumerics::AllOrders]{All\_Orders} > 0$, then
that value will have first priority.
 
\NOTE
If $\inp[nmlNumerics::SolutionOrder]{Solution\_Order} > 
\inptext{Quadrature\_Order}$, then \inptext{Quadrature\_Order}
will be set to \inptext{Solution\_Order}.
 
\NOTE
This quadrature option is only used when
  \inptext{Quadrature\_Order} $>$ \inptext{Solution\_Order}
or
  \inp[nmlNumerics::LocQuadraturePts]{Loc\_Quadrature\_Pts} $\neq$
  \inp[nmlNumerics::LocSolutionPts]{Loc\_Solution\_Pts},
otherwise the quadrature is evaluated directly at the solution points
without the need for projection.
% 
In all cases,
$\inptext{Quadrature\_Order} \geq \inptext{Solution\_Order}$
because using a lower quadrature order provides poor results.
% 
Unlike over-integration by residual projection,
the order-of-accuracy of the simulation remains \p[+1]{S}.

\NOTE
This quadrature option has not been working since the viscous and
unstructured capability was added to the code.
It should be fixed in an upcoming version.

\IDEA{
Perhaps \sref{SEC:Numerics} needs a detailed discussion of the
difference between over-integration by residual projection
and the over-integration by over-sampling at the quadrature points
used here?
}
\end{nmldescbody}

%==========================================================================

\begin{nmldescitem}
Error_Order = -3
\end{nmldescitem}
\inp*[nmlNumerics::ErrorOrder]{}[Error\_Order][]

\begin{nmldescbody}
Specifies the degree of the polynomial space to use when computing
integrated error quantities, \forexample the $L^2$-norm of the residual.
%
If
$\inptext{Error\_Order} < \inptext{Solution\_Order}$,
this will be set to the same value as \inptext{Solution\_Order}.

\NOTE
This input is only used when
  $\inptext{Error\_Order} > \inptext{Solution\_Order}$
or
  $\inp[nmlNumerics::LocErrorPts]{Loc\_Error\_Pts} \neq
   \inp[nmlNumerics::LocSolutionPts]{Loc\_Solution\_Pts}$.

\end{nmldescbody}

%==========================================================================

\begin{nmldescitem}
All_Points = -1
\end{nmldescitem}
\inp*[nmlNumerics::AllPoints]{}[All\_Points][]

\begin{nmldescbody}
Valid values for this input are the same as those specified in
\inp[nmlNumerics::LocSolutionPts]{Loc\_Solution\_Pts}.
%
If set to one of those values, this will set
\inp[nmlNumerics::LocSolutionPts]{Loc\_Solution\_Pts},
\inp[nmlNumerics::LocFluxPts]{Loc\_Flux\_Pts}, and
\inp[nmlNumerics::LocQuadraturePts]{Loc\_Quadrature\_Pts}
all to this value;
otherwise, this input will be ignored and the three mentioned input parameters
are unchanged.
%

\NOTE
This will override the values given for the three mentioned input
parameters set to a valid value.
%
\end{nmldescbody}

%==========================================================================

\begin{nmldescitem}
Loc_Solution_Pts = 1
\end{nmldescitem}
\inp*[nmlNumerics::LocSolutionPts]{}[Loc\_Solution\_Pts][]

\begin{nmldescbody}
Specifies the location of the 1D \key{solution points}[points!solution]
that are used in tensor products to create higher dimensional sets of
solution points (used for edge, quad, pyramid, prism, and hex cell geometries):
\\
\begin{tabular}{ @{\qquad} r @{ = } p{0.85\linewidth} @{} }
  0 & \key{Equidistant points}[points!equidistant] (not recommended) \\
  1 & \key{Legendre-Gauss (Gauss) points}[points!Legendre-Gauss (Gauss)] (recommended) \\ 
  2 & \key{Legendre-Gauss-Lobatto (Lobatto) points}[points!Legendre-Gauss-Lobatto (Lobatto)]
\end{tabular}
\\
See \fref{FIG:1D_Points} in \sref{SEC:Intro:Solution_Points}
to better understand the differences in these point locations.

\NOTE

For structured elements (quadrilaterals and hexahedra),
a tensor product of these 1D solution points is used.

For pyramid elements, a special tensor product of these solution points
is used. The 1D solution point locations are used to space the points
vertically from the pyramid base to the peak. At each of these vertical
locations, a 2D tensor product of solution points is used, with the base
being a tensor product of degree \PO, the next vertical location \PO[-1],
with the peak being \PO[0] or a single point.

For prism elements, a tensor product of these 1D solution points and the
specified triangle points is used. 
\end{nmldescbody}

%==========================================================================

\begin{nmldescitem}
Loc_Flux_Pts = 1
\end{nmldescitem}
\inp*[nmlNumerics::LocFluxPts]{}[Loc\_Flux\_Pts][]

\begin{nmldescbody}
Specifies the location of the flux points for edge and quadrilateral faces.
\\
\begin{tabular}{ @{\qquad} r @{ = } p{0.85\linewidth} @{} }
  0 & \key{Equidistant points}[points!equidistant] (not recommended) \\
  1 & \key{Legendre-Gauss (Gauss) points}[points!Legendre-Gauss (Gauss)] (recommended) \\ 
  2 & \key{Legendre-Gauss-Lobatto (Lobatto) points}[points!Legendre-Gauss-Lobatto (Lobatto)]
\end{tabular}
\\
For triangle faces, \inp[nmlNumerics::LocTrianglePts]{Loc\_Triangle\_Pts}
is used for the flux points.

%It is generally recommended to set the location of the flux points to the same
%as the solution points.
%
%If using Legendre-Gauss-Lobatto solution points, using Legendre-Gauss flux
%points could potentially reduce aliasing instabilities at interfaces.

\NOTE
Using different types of points for the solution and flux points is untested;
therefore, it is recommended to use the input \inptext{All\_Points} to use the
same set of points for solution, flux, and quadrature points.
%
\end{nmldescbody}

%==========================================================================

\begin{nmldescitem}
Loc_Quadrature_Pts = 1
\end{nmldescitem}
\inp*[nmlNumerics::LocQuadraturePts]{}[Loc\_Quadrature\_Pts][]

\begin{nmldescbody}
Specifies the location of the \key{quadrature points}[points!quadrature]
used for over-integration by over-sampling at the quadrature points.
This input is treated identically to
\inp[nmlNumerics::LocSolutionPts]{Loc\_Solution\_Pts} except it uses
\inp[nmlNumerics::QuadratureOrder]{Quadrature\_Order}
to determine the number of quadrature points for each geometry type.
\end{nmldescbody}

%==========================================================================

\begin{nmldescitem}
Loc_Error_Pts = 1
\end{nmldescitem}
\inp*[nmlNumerics::LocErrorPts]{}[Loc\_Error\_Pts][]

\begin{nmldescbody}
Specifies the location of the error points used for computing integrated
error quantities. This input is treated identically to
\inp[nmlNumerics::LocSolutionPts]{Loc\_Solution\_Pts} except it uses
\inp[nmlNumerics::ErrorOrder]{Error\_Order}
to determine the number of error points for each geometry type.
\end{nmldescbody}

%==========================================================================

\begin{nmldescitem}
Loc_Triangle_Pts = 1
\end{nmldescitem}
\inp*[nmlNumerics::LocTrianglePts]{}[Loc\_Triangle\_Pts][]

\begin{nmldescbody}
Specifies the location of the solution points used for triangles.
\\
\begin{tabular}{ @{\qquad} r @{ = } p{0.85\linewidth} @{} }
  -2 & Hesthaven-Warburton $\alpha$-optimized points\cite{Hesthaven_2010} \\
   0 & Equally spaced points (not recommended) \\
   1 & Williams-Shunn points\cite{Williams_2013} (recommended) \\
   2 & Barycentric Lobatto points
\end{tabular}
\end{nmldescbody}

%==========================================================================

\FINALFIX{\newpage}
\begin{nmldescitem}
Loc_Tetrahedra_Pts = 1
\end{nmldescitem}
\inp*[nmlNumerics::LocTetrahedraPts]{}[Loc\_Tetrahedra\_Pts][]

\begin{nmldescbody}
Specifies the location of the solution points used for tetrahedra.
\\
\begin{tabular}{ @{\qquad} r @{ = } p{0.85\linewidth} @{} }
  -2 & Hesthaven-Warburton $\alpha$-optimized points\cite{Hesthaven_2010} \\
   0 & Equally spaced points (not recommended)\\
   1 & Shunn-Ham points\cite{Shunn_2012} (recommended) \\
   2 & Barycentric Lobatto points \\
\end{tabular}
\end{nmldescbody}

%==========================================================================

\begin{nmldescitem}
Invs_Flux_Method = 1
\end{nmldescitem}
\inp*[nmlNumerics::InvsFluxMethod]{}[Invs\_Flux\_Method][]

\begin{nmldescbody}
Specifies the approximate Riemann solver used to compute the common inviscid
fluxes at the interfaces.
\\
\begin{tabular}{ @{\qquad} r @{ = } p{0.85\linewidth} @{} }
    1 & Roe flux\cite{Roe_1981} with entropy fix\cite{Huynh_1995} \\
    2 & HLLC flux\cite{Toro_1994} \\
    3 & LDFSS flux\cite{Edwards_1997} \\
    4 & Rotated-Roe-HLL flux\cite{Nishikawa_2008} (2D and 3D only) \\
    5 & AUSM+ flux\cite{Liou_1996} \\
    6 & Rusanov flux\cite{Rusanov_1961}
\end{tabular}\\
%
\NOTE
Only the Roe flux with entropy fix or the Rusanov flux are supported
when using a rotating reference frame. If any other option is selected
with a rotating reference frame, the inviscid flux method will be reset to
the Roe flux.

\end{nmldescbody}

%==========================================================================

\begin{nmldescitem}
Visc_Flux_Method = 2
\end{nmldescitem}
\inp*[nmlNumerics::ViscFluxMethod]{}[Visc\_Flux\_Method][]

\begin{nmldescbody}
Specifies the method used to compute the common viscous fluxes at the
interfaces.
\\
\begin{tabular}{ @{\qquad} r @{ = } p{0.85\linewidth} @{} }
  1 & Bassi-Rebay 1 (BR1)\cite{Bassi_1997} (high chances of stability issues) \\
  2 & Bassi-Rebay 2 (BR2)\cite{Bassi_2005} \\
  3 & Local DG (LDG)\cite{Cockburn_1998} (smaller CFL limit than BR2)
\end{tabular}
\end{nmldescbody}

%==========================================================================

\begin{nmldescitem}
Visc_Variables_To_Ave = 1
\end{nmldescitem}
\inp*[nmlNumerics::ViscVariablesToAve]{}[Visc\_Variables\_To\_Ave][]

\begin{nmldescbody}
Determines which variables are averaged to the interfaces when computing the
common viscous fluxes.
\\
\begin{tabular}{ @{\qquad} r @{ = } p{0.85\linewidth} @{} }
  1 & Average the interface fluxes. Compute the common viscous flux by
      averaging the left and right viscous interface fluxes which are computed
      using the left and right interface solutions and gradients. \\
  2 & Average the interface solutions and gradients. Compute common solutions
      and gradients by averaging the left and right interface solutions and
      gradients, and then use the common interface values to compute the common
      viscous flux.
\end{tabular}
\end{nmldescbody}

%==========================================================================

%##%\FINALFIX{\newpage}
\begin{nmldescitem}
Combine_Common_Flux_Projections = 0
\end{nmldescitem}
\inp*[nmlNumerics::CombineCommonFluxProjections]{}[Combine\_Common\_Flux\_Projections][]

\begin{nmldescbody}
% Specifies if the common flux projections are combined when using
Specifies if the \key[key:CommonFluxValue]{common flux}[common flux value] projections
are combined when using over-integration by over-sampling at the quadrature points.
\\
\begin{tabular}{ @{\qquad} r @{ = } p{0.85\linewidth} @{} }
  0 & All projections of the over-sampled common fluxes are performed
      separately. \\
  1 & The projections of the over-sampled interior-face common fluxes are
      performed separately from the projection of the over-sampled boundary-face
      common fluxes.\\
  2 & All of the over-sampled common fluxes are computed and combined before
      projecting.
\end{tabular}
\end{nmldescbody}

%==========================================================================

\FINALFIX{\newpage}
\begin{nmldescitem}
Correction_Function = 1
\end{nmldescitem}
\inp*[nmlNumerics::CorrectionFunction]{}[Correction\_Function][]

\begin{nmldescbody}
Specifies the \key[key:CorrectionFunction]{correction function}[=]
(\sref{SEC:Intro:FR:Correction_Functions})
to use when computing the
\key[key:ContinuousFluxFunction]{continuous flux function}[=].
\\
\begin{tabular}{ @{\qquad} r @{ = } p{0.85\linewidth} @{} }
  1 & \textsl{g}\textsubscript{DG} correction function (recovers nodal DG
    method) \\
  2 & \textsl{g}\textsubscript{2} correction function (only works for 1D)\\
  3 & \textsl{g}\textsubscript{Ga} correction function (only works for 1D)
\end{tabular}
\end{nmldescbody}

%==========================================================================

